% =========================
\section{Motivación}
% =========================
\begin{frame}{Motivación}
	
	\begin{block}{Situación inicial}
		En un cultivo experimental:
		
		\begin{itemize}
			\item Se aplica \(4x^2\) litros de agua (efecto acumulado de riego).
			\pause
			\item Se pierde \(5xy^3\) litros por exceso de fertilizante.
			\pause
			\item Se generan \(3x^3y\) litros equivalentes por crecimiento vegetal.
			\pause
			\item Se descuenta \(2x\) litros por evaporación base.
		\end{itemize}
	\end{block}
	
	\pause
	
	\textbf{Modelo algebraico:}
	
	\[
	4x^2 - 5xy^3 + 3x^3y - 2x
	\]
	
	\pause
	
	\IdeaClave{Los modelos reales pueden incluir múltiples variables y exponentes mayores a 1.}
	
\end{frame}

% =========================
\begin{frame}{Ejemplo más desafiante}
	
	\Agronomia{
		En una plantación intensiva:
		
		\begin{itemize}
			\item Producción hídrica: \(6x^2y\)
			\pause
			\item Pérdida por saturación: \( -3x^2y \)
			\pause
			\item Crecimiento adicional: \( 2x^3y^2 \)
			\pause
			\item Reducción base: \( -x^2 \)
		\end{itemize}
	}
	
	\pause
	
	\textbf{Expresión:}
	\[
	6x^2y - 3x^2y + 2x^3y^2 - x^2
	\]
	
	\pause
	
	\textbf{Simplificación parcial (solo semejantes):}
	\[
	6x^2y - 3x^2y = 3x^2y
	\]
	
	\pause
	
	\[
	3x^2y + 2x^3y^2 - x^2
	\]
	
	\pause
	
	\IdeaClave{Solo se simplifican términos con exactamente la misma parte literal.}
	
\end{frame}

% =========================
\section{Términos algebraicos avanzados}
% =========================
\begin{frame}{Términos algebraicos multivariables}
	
	\begin{block}{Definición extendida}
		Un término algebraico puede incluir:
		
		\begin{itemize}
			\item múltiples variables
			\pause
			\item exponentes mayores que 1
			\pause
			\item interacciones entre variables
		\end{itemize}
	\end{block}
	
	\pause
	
	\[
	3x^2y^3
	\]
	
	\pause
	
	\Agronomia{
		Representa un efecto combinado entre riego (\(x\)) y fertilización (\(y\)) de alta intensidad.
	}
	
\end{frame}

% =========================
\begin{frame}{Ejemplos en agronomía}
	
	\begin{columns}
		\begin{column}{0.48\textwidth}
			
			\begin{block}{Ejemplo 1}
				\[
				5x^2
				\]
				\pause
				Crecimiento del cultivo en función del riego.
			\end{block}
			
			\pause
			
			\begin{block}{Ejemplo 2}
				\[
				-3xy^3
				\]
				\pause
				Pérdida por sobrefertilización.
			\end{block}
			
		\end{column}
		
		\begin{column}{0.48\textwidth}
			
			\begin{exampleblock}{Ejemplo complejo}
				\[
				4x^3y^2
				\]
				\pause
				
				\begin{itemize}
					\item Coeficiente: 4
					\pause
					\item Variables: \(x,y\)
					\pause
					\item Exponentes: 3 y 2
				\end{itemize}
				
			\end{exampleblock}
			
		\end{column}
	\end{columns}
	
\end{frame}

% =========================
\section{Expresiones algebraicas}
% =========================
\begin{frame}{Expresión algebraica multivariable}
	
	\begin{block}{Definición}
		Es la combinación de términos con varias variables y distintos exponentes.
	\end{block}
	
	\pause
	
	\[
	4x^2 - 5xy^3 + 3x^3y - 2x
	\]
	
	\pause
	
	\Agronomia{
		Modelo de rendimiento considerando interacción agua-fertilizante.
	}
	
\end{frame}

% =========================
\begin{frame}{Ejemplo resuelto}
	
	\Agronomia{
		En un cultivo de maíz:
		
		\begin{itemize}
			\item Aporte hídrico: \(5x^2\)
			\pause
			\item Pérdida: \( -2x^2 \)
			\pause
			\item Interacción con fertilizante: \( 3xy^2 \)
			\pause
			\item Ajuste técnico: \( +4xy^2 \)
		\end{itemize}
	}
	
	\pause
	
	\textbf{Expresión:}
	\[
	5x^2 -2x^2 + 3xy^2 + 4xy^2
	\]
	
	\pause
	
	\[
	3x^2 + 7xy^2
	\]
	
	\pause
	
	\IdeaClave{Se agrupan términos semejantes respetando variables y exponentes.}
	
\end{frame}

% =========================
\section{Términos semejantes}
% =========================
\begin{frame}{Términos semejantes avanzados}
	
	\begin{block}{Definición}
		Son semejantes si tienen:
		
		\begin{itemize}
			\item mismas variables
			\pause
			\item mismos exponentes
		\end{itemize}
	\end{block}
	
	\pause
	
	\[
	2x^2y^3 + 5x^2y^3 = 7x^2y^3
	\]
	
\end{frame}

% =========================
\begin{frame}{Ejemplo analizado}
	
	\[
	4x^2,\quad -3x^2,\quad 5xy,\quad 2xy^2
	\]
	
	\pause
	
	\begin{itemize}
		\item \(4x^2\) y \(-3x^2\) son semejantes
		\pause
		\item \(5xy\) y \(2xy^2\) NO son semejantes
	\end{itemize}
	
	\pause
	
	\[
	4x^2 - 3x^2 = x^2
	\]
	
	\pause
	
	\IdeaClave{Cambiar un exponente cambia completamente el tipo de término.}
	
\end{frame}

% =========================
\section{Actividades}
% =========================
\begin{frame}{Actividad 1}
	
	\textbf{Simplifique:}
	
	\[
	3x^2y^3 - 5x^2y^3 + 2x^3y - x^3y
	\]
	
	\pause
	
	\[
	-2x^2y^3 + x^3y
	\]
	
\end{frame}

% =========================
\begin{frame}{Actividad 2}
	
	\textbf{Simplifique:}
	
	\[
	6xy^2 + 4xy^2 - 3xy^2 + 2x^2y
	\]
	
	\pause
	
	\[
	7xy^2 + 2x^2y
	\]
	
\end{frame}

% =========================
\begin{frame}{Actividad 3}
	
	\textbf{Evaluar si } \(x=2, y=1\):
	
	\[
	4x^2 - 5xy^3 + 3x^3y - 2x
	\]
	
	\pause
	
	\[
	4(4) - 5(2)(1) + 3(8)(1) - 4
	\]
	
	\pause
	
	\[
	16 -10 +24 -4 = 26
	\]
	
\end{frame}

% =========================
\begin{frame}{Actividad 4 (análisis)}
	
	\Agronomia{
		Un técnico afirma que:
		\[
		4x^2y + 3xy^2 = 7x^2y^2
		\]
	}
	
	\pause
	
	\textbf{¿Es correcto?}
	
	\pause
	
	No.
	
	\pause
	
	\textbf{Razón:}
	
	\begin{itemize}
		\item Cambian exponentes
		\item Cambia estructura de variables
	\end{itemize}
	
	\pause
	
	\IdeaClave{No se pueden combinar términos no semejantes.}
	
\end{frame}

% =========================
\section{Patrones}
% =========================
\begin{frame}{Detección de patrones}
	
	\[
	x^2,\ 2x^2,\ 4x^2,\ 8x^2
	\]
	
	\pause
	
	\begin{itemize}
		\item Crece multiplicando por 2
		\pause
		\item Mantiene misma estructura algebraica
	\end{itemize}
	
\end{frame}

% =========================
\begin{frame}{Inconsistencia}
	
	\[
	x^2,\ 2x^2,\ 5x^2,\ 8x^2
	\]
	
	\pause
	
	\begin{itemize}
		\item El término \(5x^2\) rompe el patrón
	\end{itemize}
	
\end{frame}